\chapter[ContentManager and Asset Resolution]{ContentManager and Asset \mbox{Resolution}}
\label{ch:contentmanager-resolution}
\label{ch:content-and-assets}

\cnaclass{ContentManager} maps a logical name and requested C\texttt{++} type to an XNB object,
a CNJ object, or a native loose asset. Its most important rule is selection, not decoding: the
first existing candidate is terminal even if the selected reader later rejects it. This chapter
owns path choice, caching, and manager lifetime. Chapters~\ref{ch:xnb-container} and
\ref{ch:xnb-type-readers} own XNB; Chapter~\ref{ch:cnj-format} owns CNJ.

\begin{figure}[tbp]
\centering
\begin{figurealt}{ContentManager Load first checks its cache. On a miss it probes the logical base plus dot xnb. If absent, it obtains the loose type reader, checks an existing literal base, then base plus dot cnj, then the reader's ordered native extensions. The first existing candidate is terminal; reader failure does not continue to a lower-priority candidate. A successful result enters the relevant cache.}
\resizebox{\textwidth}{!}{%
\begin{tikzpicture}[node distance=8mm and 9mm]
  \node[cna public] (load) {\texttt{Load<T>(name)}};
  \node[cna internal, right=of load] (cache) {type + normalized-name\\cache lookup};
  \node[cna internal, right=of cache] (xnb) {\texttt{base + .xnb} exists?};
  \node[cna internal, right=of xnb] (literal) {literal \texttt{base} exists?};
  \node[cna internal, right=of literal] (cnj) {\texttt{base + .cnj} exists?};
  \node[cna internal, right=of cnj] (ext) {first existing reader\\extension};
  \node[cna warning, below=of literal] (terminal) {selected candidate is terminal\\decode failure does not retry};
  \node[cna public, below=of ext] (result) {decoded asset / exception};
  \draw[cna flow] (load) -- (cache);
  \draw[cna flow] (cache) -- node[cna label]{miss} (xnb);
  \draw[cna flow] (xnb) -- node[cna label]{absent} (literal);
  \draw[cna flow] (literal) -- node[cna label]{absent} (cnj);
  \draw[cna flow] (cnj) -- node[cna label]{absent} (ext);
  \draw[cna flow] (xnb) |- (terminal);
  \draw[cna flow] (literal) -- (terminal);
  \draw[cna flow] (cnj) |- (terminal);
  \draw[cna flow] (ext) -- (result);
  \draw[cna flow] (terminal) -- (result);
  \draw[cna optional] (result) -| node[cna label,near start]{success caches} (cache.south);
\end{tikzpicture}}
\end{figurealt}
\caption{ContentManager candidate order at \CnaRevisionShort. Existence selects a path; it does
not prove that the selected reader can decode or use the asset.}
\label{fig:content-resolution}
\end{figure}



\section{The resolution algorithm}

On a cache miss, let \texttt{base = BuildAssetPath(assetName)}. The generic
\texttt{Load<T>} path and the Texture2D, SoundEffect, and TextureCube specializations follow this
order:

\begin{lstlisting}
if (exists(base + ".xnb"))
    return LoadXnbAsset<T>(base + ".xnb");

reader = registered LooseFileContentTypeReader<T>;
if (exists(base))
    return reader.Read(base, content);
if (exists(base + ".cnj"))
    return reader.Read(base + ".cnj", content);
for (const auto& extension : reader.GetExtensions())
    if (exists(base + extension))
        return reader.Read(base + extension, content);
return reader.Read(base, content); // preserve the reader's diagnostic
\end{lstlisting}

\begin{contractnote}
Existence selects a candidate; successful decoding does not. A malformed XNB, invalid CNJ,
corrupt literal file, directory, or rejected first native extension stops the load. CNA does not
retry a lower-priority sibling.
\end{contractnote}

The XNB probe appends \texttt{.xnb} to the complete name. Therefore
\texttt{Load<Texture2D>("art/logo.png")} first tests \texttt{art/logo.png.xnb}, then the literal
PNG. To parse \texttt{art/logo.xnb} as XNB, pass \texttt{"art/logo"}. Passing the suffix itself
tests \texttt{logo.xnb.xnb}; if absent, the loose reader receives the original XNB file.
Dots in logical names have no special meaning.

A literal path wins over its sidecar. With \texttt{hero.png} and \texttt{hero.cnj} present:

\begin{lstlisting}
auto native  = content.Load<Texture2D>("sprites/hero.png");
auto sidecar = content.Load<Texture2D>("sprites/hero");
\end{lstlisting}

The first call selects the literal PNG after the extra XNB probe. The second selects
\texttt{hero.cnj} before the native-extension list. A CNJ \texttt{sourceFile} then starts a
checked nested load; it is not a direct decoder bypass.

\subsection{Native extension order}

\begin{center}
\small
\begin{tabular}{@{}p{0.25\textwidth}p{0.62\textwidth}@{}}
\toprule
\textbf{Requested type} & \textbf{Ordered native candidates after literal and CNJ} \\
\midrule
\cnaclass{Texture2D} & \texttt{.png}, \texttt{.jpg}, \texttt{.jpeg}, \texttt{.bmp},
  \texttt{.gif}, \texttt{.tga}, \texttt{.tif}, \texttt{.tiff}, \texttt{.qoi} \\
\cnaclass{TextureCube} & \texttt{.dds} \\
\cnaclass{SoundEffect} & \texttt{.wav} \\
\cnaclass{Song} & \texttt{.mp3}, \texttt{.ogg}, \texttt{.wav}, \texttt{.flac},
  \texttt{.opus}, \texttt{.aac}, \texttt{.wma} \\
\cnaclass{Video} & \texttt{.mp4}, \texttt{.ogv}, \texttt{.webm}, \texttt{.mkv},
  \texttt{.avi}, \texttt{.mov}; loose reader excluded on Emscripten, Android, and MinGW \\
\cnaclass{Model} & \texttt{.cnj}, \texttt{.gltf}, \texttt{.glb}; the first entry is redundant
  with the manager-level CNJ probe \\
\cnaclass{SpriteFont}, \cnaclass{Effect}, \cnaclass{AnimationClipEXT}, \cnaclass{Curve} &
  self-contained \texttt{.cnj} \\
\bottomrule
\end{tabular}
\end{center}

The global CNJ probe runs even when a reader does not advertise or understand CNJ. In particular,
the current loose \cnaclass{Song} and \cnaclass{Video} readers treat the selected path as media.
An extensionless name with a same-named \texttt{.cnj} therefore feeds JSON to the media object
instead of falling through to the native audio/video file. Pass the explicit media filename for
those types; a generic media-sidecar contract does not exist at the pin.

\section{RootDirectory and containment}

\texttt{RootDirectory} is a base path, not a sandbox. \texttt{BuildAssetPath} joins it with the
logical name but does not reject absolute paths, canonicalize dot segments, or check symlinks:

\begin{lstlisting}
root = "Content"
"sprites/logo"   -> "Content/sprites/logo"
"../shared/logo" -> "Content/../shared/logo"
"/srv/logo"      -> "/srv/logo"
\end{lstlisting}

Top-level names must therefore be trusted or validated by the application. Internal references
have narrower rules:

\begin{itemize}
  \item CNJ \texttt{sourceFile} rejects absolute paths, traversal and symlink escape, another
        CNJ, and implicit sidecar cycles.
  \item \cnaclass{ContentReader::ReadExternalReference<T>} rejects absolute, drive, UNC, and
        normalized above-root paths.
  \item XNB Song and Video readers contain companion files within the content root. For an
        explicitly loaded external XNB, the XNB's own directory becomes the authorized bundle
        root.
\end{itemize}

These checks compare canonicalized paths before the eventual open; they are not a race-proof
open-by-handle policy. Chapter~\ref{ch:content-robustness} records the TOCTOU boundary.

\subsection{Cache identity is textual}

\texttt{NormalizeKey} converts backslashes to slashes and lowercases bytes. It does not include
the root, resolve dot segments, or canonicalize a file. Consequences include:

\begin{itemize}
  \item \texttt{Sprites\textbackslash Logo} and \texttt{sprites/logo} share a cache key;
  \item \texttt{textures/../logo} and \texttt{logo} may name one file but occupy two entries;
  \item two case-distinct files on a case-sensitive filesystem collapse to one logical key;
  \item changing \texttt{RootDirectory} does not invalidate an already cached logical name.
\end{itemize}

Call \texttt{Unload} before intentionally reusing names under a new root.

\section{Direct and managed loading}
\label{sec:direct-loading}

A direct texture constructor with a concrete path remains useful for a small program:

\begin{lstlisting}
Texture2D logo("assets/logo.png", getGraphicsDeviceProperty());
\end{lstlisting}

\cnaclass{Game} instead exposes its configured manager in \texttt{LoadContent}:

\begin{lstlisting}
ContentManager& content = getContentProperty();
content.setRootDirectoryProperty("Content");
auto logo = content.Load<Texture2D>("sprites/logo");
auto font = content.Load<SpriteFont>("fonts/hud");
\end{lstlisting}

\texttt{RegisterTypeReader<T>} installs one loose reader for a C\texttt{++} type.
\texttt{RegisterCnjLoader<T>} installs named CNJ factories that produce a type without a
pre-existing reader. Both are CNAEXT surfaces; complete registration before loading begins.

\section{Caching and ownership}

Successful generic loads enter a strong cache keyed by requested type and normalized logical
name. The write occurs only after a reader returns, so a failed top-level load can be retried.
Specializations differ:

\begin{center}
\small
\begin{tabular}{@{}p{0.24\textwidth}p{0.31\textwidth}p{0.34\textwidth}@{}}
\toprule
\textbf{Type} & \textbf{Manager state} & \textbf{Ownership consequence} \\
\midrule
Ordinary copyable \texttt{T} & Strong \texttt{std::any} value & Repeated loads return copies of
  the cached value; caller-held shared handles can outlive eviction. \\
\cnaclass{Texture2D} & Weak renderer and optional pixel shadow & A live returned texture keeps
  its renderer alive; after eviction, a new load can construct a new renderer. \\
\cnaclass{SoundEffect}, \cnaclass{TextureCube} & No manager cache & Move-only values are decoded
  independently on each call. \\
\bottomrule
\end{tabular}
\end{center}

\texttt{Unload} clears only the generic and texture caches. It does not dispose caller-held
values, clear readers, reset the root/device/provider, or refresh the manifest.
\texttt{Dispose} calls \texttt{Unload} once and marks the manager; later loads throw
\texttt{std::runtime\_error}. The default C\texttt{++} destructor does not call that public
method.

Nested reader activity is not transactional. A reader can cache dependencies and then fail,
leaving those successful nested entries. There is no provisional key or cycle detector, so a
recursive load of the same uncached type/name can recurse until stack exhaustion.

\subsection{Game content assignment copies state}

\cnaclass{Game}\texttt{::setContentProperty(const ContentManager\&)} performs memberwise assignment into
the existing manager. It copies caches, readers, manifest, disposed state, and raw service/device
pointers; it does not replace a managed-object reference as FNA does. Assigning a standalone
manager can overwrite Game's valid graphics-device pointer with null. Configure the existing
manager in place when only the root should change:

\begin{lstlisting}
auto& content = getContentProperty();
content.setRootDirectoryProperty("OtherContent");
\end{lstlisting}

\section{Failure and concurrency boundaries}

\texttt{ContentLoadException} covers many validated content errors, including malformed XNB/CNJ
data, absent reader registration, bad reader indices, and unsupported reader versions. It is not
a universal wrapper. Depending on the selected path, callers can also see
\texttt{std::runtime\_error}, \cnaclass{FileNotFoundException},
\cnaclass{EndOfStreamException}, \texttt{std::bad\_any\_cast}, or an exception from custom
reader code. Empty names are not rejected up front.

\begin{lstlisting}
try {
    auto level = content.Load<GameLevelData>("levels/intro");
    StartLevel(level);
} catch (const ContentLoadException& error) {
    ShowContentDiagnostic(error.what());
} catch (const std::exception& error) {
    ReportAssetFailure(error.what());
}
\end{lstlisting}

The manager's maps, manifest vector, strings, pointers, and disposed flag are unsynchronized.
Concurrent load/unload/registration/root/manifest operations can form C\texttt{++} data races.
Single-threaded reentrancy also has hazards: self-loading cycles have no sentinel, and a reader
that replaces its own registration can destroy the active object. Register and configure first,
serialize manager mutation through one owner, and publish completed assets to workers.

\section{Provider, ResourceContentManager, and TitleContainer}

The \texttt{IServiceProvider*} constructors retain a raw pointer, but the pinned
\cnaclass{ContentManager} never calls \texttt{GetService}. A GPU-backed standalone manager needs
an explicit device:

\begin{lstlisting}
ContentManager assets(&services, "Content");
assets.setGraphicsDevice(device);
auto logo = assets.Load<Texture2D>("sprites/logo");
\end{lstlisting}

\cnaclass{Game} performs that association for its own manager. Both pointers remain borrowed.

\cnaclass{ResourceContentManager} is a surface stub. Its protected \texttt{OpenStream} ignores
the asset name and throws \texttt{std::runtime\_error}; inherited \texttt{Load<T>} does not call
that virtual method. The class therefore loads ordinary files only through its base behavior and
does not provide embedded resources.

\subsection{TitleContainer}
\label{sec:title-container-contract}

\cnaclass{TitleContainer::OpenStream} is an independent direct-file helper. It converts
backslashes, resolves relative names against the process-global \cnaclass{TitleLocation}, and
returns a read stream. ContentManager does not call it. The CNAEXT title-location setter is
mutable global configuration and is unsynchronized.

Resolution is lexical rather than contained: parent segments, absolute paths, and symlinks can
leave the title directory. Android adds an SDL asset fallback and copies the loaded bytes into a
\cnaclass{MemoryStream}. The helper is suitable for application-controlled paths, not untrusted
filenames or content-root redirection.

\section{Content manifest}

\texttt{GetContentManifest} lazily scans the root and returns a reference to manager-owned rows.
\texttt{RefreshContentManifest} rebuilds them. A row records a logical path, XNB/CNJ presence,
native extensions, and---for readable uncompressed XNBs---reader names.
\texttt{GetXnbReaderUsageSummary} aggregates those names and reports whether the global registry
currently contains each key.

The manifest is diagnostic only:

\begin{itemize}
  \item Load performs fresh filesystem probes and never consults it.
  \item Presence does not establish a valid file, compatible reader, device, or loadable graph.
  \item Compressed or malformed XNBs can appear with an empty reader-name list.
  \item Directory and hash-map iteration make row and summary order unstable.
  \item Missing or inaccessible trees can yield an empty or partial snapshot without a public
        scan-status result.
  \item Refresh can invalidate held element references and iterators; copy before sorting or
        retaining the data.
\end{itemize}

\texttt{Unload}, disposal, and root changes do not automatically refresh this separate snapshot.

\section{Audio and video readers}
\label{sec:xnb-audio}

Loose readers and XNB readers solve different problems. Loose \cnaclass{Song} and
\cnaclass{Video} simply wrap the selected media path (Video also requires the manager device).
XNB readers decode the XNA-specific body and resolve an external companion.

\subsection{SoundEffect}

The XNB \cnaclass{SoundEffectReader} handles native PCM16 directly. PCM8, IEEE float, IMA-ADPCM,
and MS-ADPCM are wrapped as an in-memory WAV and sent through SDL's decoder; XMA2 is rejected.
MS-ADPCM content without its coefficient extension receives the standard coefficient table and
a block-derived sample count. The full wire contract and loop-boundary defect are in
\S\ref{sec:xnb-soundeffect-reader-contract}.

\subsection{Song}
\label{sec:xnb-song-reader}

An XNB Song body contains a reference string and duration in milliseconds. The reader resolves
the reference within the authorized root, strips the final four characters when long enough,
and probes the extensionless stem, \texttt{.ogg}, \texttt{.oga}, then \texttt{.qoa}. If none
exists, it restores the serialized path. This preserves the desktop substitution convention for
XNA \texttt{.wma} placeholders; it is not general extension parsing.

Construction checks that the selected path exists, so a well-formed XNB with no companion fails
during \texttt{Load<Song>} with \cnaclass{FileNotFoundException}. Decoding remains deferred until
playback. The retained MonoGame fixture establishes XNB dispatch, sibling lookup, duration, and
the \texttt{.ogg} route; it does not establish playback or the other substitutions.

\subsection{Video}
\label{sec:xnb-video-reader}

An XNB Video body contains a reference, duration, width, height, frame rate, and soundtrack enum.
The reader contains the reference, probes the extensionless stem, \texttt{.ogv}, then
\texttt{.ogg}, and stores the serialized metadata in a \cnaclass{Video}. It does not open or
validate the media during content loading.

\cnaclass{VideoPlayer::Play} performs the FFmpeg open on supported native builds and compares
stream dimensions and frame rate with the serialized metadata. Missing or unusable media can
leave CNA stopped without the FNA-style file exception. The strongest test uses a test-built XNB
and CNA's FFmpeg fixture; no external-producer Video XNB is retained. On platforms where the
FFmpeg-backed implementation is excluded, declarations alone do not establish a usable route.

\section{Porting implications}

Use an extensionless logical name when XNB/CNJ precedence is desired; use a concrete literal
filename when the native file itself must win. Treat the root and TitleContainer as trusted-path
conveniences, not sandboxes. Register readers and attach a graphics device before loading, keep
manager mutation serialized, and call \texttt{Unload} when deliberately changing the cache
namespace.

Existing XNA/FNA games can reuse compiled assets only for reader families supported at this pin.
Arbitrary Effect bytecode remains a separate unsupported contract
(Chapter~\ref{ch:shader-fx-gap}); CNJ or a native conversion is appropriate where the XNB reader
or runtime capability is absent.

\section{Summary}

ContentManager's contract is a deterministic existence-priority resolver with type-specific
caching. XNB, literal, CNJ, and native-extension paths are not fallback attempts after decode
failure. The distinction explains most surprising content behavior: explicit filename choice,
root changes hidden by cache, sidecars shadowing native media, and a manifest that inventories
candidates without certifying them.
